\setcounter{section}{10}

  \zwischenueberschrift{Übungsaufgaben}

\inputaufgabe
{}
{

Zeige, dass

\mavergleichskettedisp
{\vergleichskette
{M
}
{ =} { \left\{ (x,y,z,t) \in \R^4  \mid   x+x^2y+z^2+t^3 = 0         \right\}
}
{ } {
}
{ } {
}
{ } {
}
}
{}{}{}
eine
\definitionsverweis {abgeschlossene Untermannigfaltigkeit}{}{}
des $\R^4$ ist.

}
{} {}

\inputaufgabe
{}
{

Es sei
\maabbdisp {f} {\R} {\R_+
} {}
eine
\definitionsverweis {stetig differenzierbare Funktion}{}{.}
Zeige, dass die
\definitionsverweis {Rotationsfläche}{}{}
des
\definitionsverweis {Graphen}{}{}
von $f$ eine
\definitionsverweis {abgeschlossene Untermannigfaltigkeit}{}{}
des $\R^{3}$ ist.

}
{} {}

\inputaufgabe
{}
{

Zeige, dass die Menge aller reellen
\mathl{n\times n}{-}Matrizen mit
\definitionsverweis {Determinante}{}{}
$1$ eine
\mathl{(n^2-1)}{-}dimensionale
\definitionsverweis {abgeschlossene Untermannigfaltigkeit}{}{}
des $\R^{n^2}$ ist.

}
{} {}

\inputaufgabe
{}
{

Man gebe ein Beispiel einer
\definitionsverweis {abgeschlossenen Teilmenge}{}{}
\mathl{M\subseteq \R}{,} die keine
\definitionsverweis {abgeschlossene Untermannigfaltigkeit}{}{} von $\R$ ist.

}
{} {}

\inputaufgabe
{}
{

Bestimme die
\definitionsverweis {abgeschlossenen Untermannigfaltigkeiten}{}{}
von $\R$.

}
{} {}

\inputaufgabe
{}
{

Bestimme die
\definitionsverweis {abgeschlossenen Untermannigfaltigkeiten}{}{}
von $S^1$.

}
{} {}

\inputaufgabe
{}
{

Es seien
\mathkor {} {M} {und} {N} {}
zwei
\definitionsverweis {disjunkte}{}{}
\definitionsverweis {abgeschlossene Untermannigfaltigkeiten}{}{}
des $\R^n$. Zeige, dass deren
\definitionsverweis {Vereinigung}{}{}

\mathl{M \cup N}{} ebenfalls eine abgeschlossene Untermannigfaltigkeit ist, und dass diese Aussage ohne die Voraussetzung der Disjunktheit nicht gilt.

}
{} {}

\inputaufgabe
{}
{

Es sei
\maabb {f} {\R^n } {\R
} {}
eine
\definitionsverweis {stetig differenzierbare Funktion}{}{.}
Zeige, dass der
\definitionsverweis {Graph}{}{}

\mavergleichskette
{\vergleichskette
{ \Gamma_f
}
{ \subseteq }{ \R^{n+1}
}
{  }{
}
{    }{
}
{   }{
}
}
{}{}{}
eine
\definitionsverweis {abgeschlossene Untermannigfaltigkeit}{}{}
des $\R^{n+1}$ ist.

}
{} {}

\inputaufgabegibtloesung
{}
{

Es sei

\mavergleichskette
{\vergleichskette
{ M
}
{ \neq }{ \emptyset
}
{  }{
}
{    }{
}
{   }{
}
}
{}{}{}
eine
\definitionsverweis {differenzierbare Mannigfaltigkeit}{}{}
der Dimension $n$. Zeige, dass es eine Kette von
\definitionsverweis {abgeschlossenen Untermannigfaltigkeiten}{}{}

\mavergleichskettedisp
{\vergleichskette
{ M_0
}
{ \subseteq} { M_1
}
{ \subseteq} { M_2
}
{ \subseteq  \ldots \subseteq} { M_{n-1}
}
{ \subseteq} { M_n
}
}
{
\vergleichskettefortsetzung
{ =} { M
}
{ } {}
{ } {}
{ } {}
}{}{}
derart gibt, dass die abgeschlossene Untermannigfaltigkeit $M_i$ die Dimension $i$ besitzt.

}
{} {}

\inputaufgabegibtloesung
{}
{

Wir betrachten die Menge

\mavergleichskettedisp
{\vergleichskette
{N
}
{ =} { \left\{ A= \begin{pmatrix}  a   &   b   \\  c  & d \end{pmatrix}   \mid   A \text{ ist nilpotent}         \right\}
}
{ \subseteq} {\R^4
}
{ } {
}
{ } {
}
}
{}{}{}
der reellen nilpotenten
$2\times 2$-\definitionsverweis {Matrizen}{}{}
sowie die Menge

\mavergleichskettedisp
{\vergleichskette
{M
}
{ =} { N \setminus \{ \begin{pmatrix}  0   &   0  \\  0  &  0 \end{pmatrix}  \}
}
{ } {
}
{ } {
}
{ } {
}
}
{}{}{.}

a) Ist $N$ zusammenhängend?

b) Zeige, dass $M$ eine
\definitionsverweis {abgeschlossene Untermannigfaltigkeit}{}{}
einer offenen Teilmenge

\mavergleichskette
{\vergleichskette
{ G
}
{ \subseteq }{ \R^4
}
{  }{
}
{    }{
}
{   }{
}
}
{}{}{}
ist.

c) Bestimme die
\definitionsverweis {Dimension}{}{}
von $M$.

d) Ist $M$ zusammenhängend?

e) Überdecke
\mathl{M}{} mit expliziten topologischen Karten.

}
{} {}

\inputaufgabe
{}
{

Es sei

\mavergleichskette
{\vergleichskette
{ M
}
{ \subseteq }{ \R^n
}
{  }{
}
{    }{
}
{   }{
}
}
{}{}{}
eine
\definitionsverweis {abgeschlossene Untermannigfaltigkeit}{}{}
und

\mavergleichskette
{\vergleichskette
{ P
}
{ \in }{ M
}
{  }{
}
{    }{
}
{   }{
}
}
{}{}{}
ein Punkt. Es sei
\mathl{T_P(i)}{} die
\definitionsverweis {Tangentialabbildung}{}{}
zur Inklusion
\maabbdisp {i} { M } {\R^n
} {}
und $\gamma$ ein
\definitionsverweis {differenzierbarer Weg}{}{}
in $M$ mit

\mavergleichskettedisp
{\vergleichskette
{ \gamma(0)
}
{ =} { P
}
{ } {
}
{ } {
}
{ } {
}
}
{}{}{.}
Zeige, dass in

\mavergleichskette
{\vergleichskette
{ \R^n
}
{ \cong }{ T_P \R^n
}
{  }{
}
{    }{
}
{   }{
}
}
{}{}{}
die Gleichheit

\mavergleichskettedisp
{\vergleichskette
{ T_P(i) ( [\gamma] )
}
{ =} { \left(  i \circ \gamma  \right) '(0)
}
{ } {
}
{ } {
}
{ } {
}
}
{}{}{}
gilt.

}
{} {}

\inputaufgabe
{}
{

Es sei

\mavergleichskette
{\vergleichskette
{ G
}
{ \subseteq }{ \R^n
}
{  }{
}
{    }{
}
{   }{
}
}
{}{}{}
\definitionsverweis {offen}{}{,}
\maabb {\varphi} { G } {\R^m
} {}
eine
\definitionsverweis {differenzierbare Abbildung}{}{}
und $M$ die
\definitionsverweis {Faser}{}{}
über

\mavergleichskette
{\vergleichskette
{ 0
}
{ \in }{ \R^m
}
{  }{
}
{    }{
}
{   }{
}
}
{}{}{.}
Es sei vorausgesetzt, dass das
\definitionsverweis {totale Differential}{}{}
in jedem Punkt dieser Faser
\definitionsverweis {surjektiv}{}{}
sei. Zeige, dass für

\mavergleichskette
{\vergleichskette
{ P
}
{ \in }{ M
}
{  }{
}
{    }{
}
{   }{
}
}
{}{}{}
der
\definitionsverweis {Tangentialraum}{}{}
im Sinne von
[[Differenzierbare Abbildung/R/Regulärer Punkt/Tangentialraum/An Faser/Definition|Definition 53.7 (Analysis (Osnabrück 2021-2023))  ]]  mit dem
\definitionsverweis {Tangentialraum}{}{}
der
\definitionsverweis {differenzierbaren Mannigfaltigkeit}{}{}
$M$ im Punkt $P$ übereinstimmt.

}
{} {}

\inputaufgabe
{}
{

Es sei $M$ eine
\definitionsverweis {differenzierbare Mannigfaltigkeit}{}{}
und
\maabbdisp {\pi} {TM} {M
} {}
das
\definitionsverweis {Tangentialbündel}{}{.}
Zeige, dass diese Projektionsabbildung
\definitionsverweis {stetig}{}{}
ist.

}
{} {}

\inputaufgabe
{}
{

Es sei $TM$ das
\definitionsverweis {Tangentialbündel}{}{}
zu einer
\definitionsverweis {differenzierbaren Mannigfaltigkeit}{}{}
$M$. Zeige, dass die Mengen
\mathl{(T(\alpha))^{-1} (V \times W)}{} zu allen Karten
\maabbdisp {\alpha} { U } { V
} {}
mit
\mathl{V,W \subseteq \R^n}{} offen eine
\definitionsverweis {Basis der Topologie}{}{}
auf dem Tangentialbündel bilden.

}
{} {}

\inputaufgabegibtloesung
{}
{

Es sei

\mavergleichskette
{\vergleichskette
{ W
}
{ \subseteq }{ \R^n
}
{  }{
}
{    }{
}
{   }{
}
}
{}{}{}
\definitionsverweis {offen}{}{,}
\maabb {h} { W } { \R
} {}
eine
\definitionsverweis {stetig differenzierbare Funktion}{}{}
und

\mavergleichskette
{\vergleichskette
{ Y
}
{ = }{ h^{-1}(c)
}
{  }{
}
{    }{
}
{   }{
}
}
{}{}{}
die
\definitionsverweis {Faser}{}{}
zu

\mavergleichskette
{\vergleichskette
{ c
}
{ \in }{ \R
}
{  }{
}
{    }{
}
{   }{
}
}
{}{}{,}
wobei $h$ in jedem Punkt von $Y$
\definitionsverweis {regulär}{}{}
sei. Realisiere das
\definitionsverweis {Tangentialbündel}{}{}
von $Y$ als eine
\definitionsverweis {abgeschlossene Untermannigfaltigkeit}{}{}
von $W \times \R^n$.

}
{} {}

\inputaufgabe
{}
{

Es sei

\mavergleichskette
{\vergleichskette
{ Y
}
{ \subseteq }{ \R^n
}
{  }{
}
{    }{
}
{   }{
}
}
{}{}{}
eine
\definitionsverweis {abgeschlossene Untermannigfaltigkeit}{}{}
und
\maabb {\varphi} { V } { U \subseteq Y
} {}
eine diffeomorphe Parametrisierung einer offenen Teilmenge

\mavergleichskette
{\vergleichskette
{ U
}
{ \subseteq }{ Y
}
{  }{
}
{    }{
}
{   }{
}
}
{}{}{}
durch eine offene Teilmenge

\mavergleichskette
{\vergleichskette
{ V
}
{ \subseteq }{ \R^d
}
{  }{
}
{    }{
}
{   }{
}
}
{}{}{,}
wobei wir $\varphi$ auch als Abbildung nach $\R^n$ auffassen. Zeige, dass ein kommutatives Diagramm

\mathdisp {\begin{matrix}  V  & \stackrel{ \partial_i }{\longrightarrow} &  TV &  \\ \!\!\!\!\! \,\, \, \varphi \downarrow & \partial_i \varphi \searrow & \downarrow T(\varphi) \!\!\!\!\!  & \\  U  & \stackrel{  }{\longrightarrow} &  TU & \!\!\!\!\!   \\ \end{matrix}} {  }

vorliegt.

}
{} {}

\inputaufgabe
{}
{

Seien
\mathkor {} {M} {und} {N} {}
\definitionsverweis {differenzierbare Mannigfaltigkeiten}{}{} und es sei
\maabb {\varphi} { M } { N
} {} eine
\definitionsverweis {differenzierbare Abbildung}{}{.} Zeige, dass die zugehörige
\definitionsverweis {Tangentialabbildung}{}{}
\maabbdisp {T(\varphi)} {TM} {TN
} {}
\definitionsverweis {stetig}{}{}
ist.

}
{} {}

\inputaufgabe
{}
{

Berechne die
\definitionsverweis {Tangentialabbildung}{}{} $T\varphi$ zu
\maabbeledisp {\varphi} {\R^3} {\R^2
} {(x,y,z)} {(x^2y-3xz^3+y^2,x \sin  y  -e^{yz})
} {} unter Verwendung der Identifizierungen
\mathkor {} {T\R^3 =\R^3 \times \R^3} {und} {T\R^2 =\R^2 \times \R^2} {.}

}
{} {}

\inputaufgabe
{}
{

Man gebe ein Beispiel einer differenzierbaren Kurve
\maabbdisp {\gamma} { [0,1[} { S^1
} {}
derart, dass der
\definitionsverweis {Grenzwert}{}{}

\mathl{\operatorname{lim}_{   t  \rightarrow  1   } \,  \gamma(t)}{} existiert, dass aber der Grenzwert
\mathl{\operatorname{lim}_{   t  \rightarrow  1   } \, ( \gamma(t), T_t (\gamma ) (1))}{} in
\mathl{TS^1}{} nicht existiert.

}
{} {}

\inputaufgabe
{}
{

Es sei

\mavergleichskette
{\vergleichskette
{ M
}
{ \subseteq }{ \R^n
}
{  }{
}
{    }{
}
{   }{
}
}
{}{}{}
eine
\definitionsverweis {abgeschlossene Untermannigfaltigkeit}{}{.}
Interpretiere die Hintereinanderschaltung

\mathdisp {TM \longrightarrow T\R^n = \R^n \times \R^n \stackrel{+}{\longrightarrow} \R^n} { . }

}
{} {}

\inputaufgabe
{}
{

Zeige, dass die Abbildung
\maabbeledisp {} {TS^1} { \R^2
} {( (a,b),t (-b,a))} { (a,b) + t (-b,a)
} {,}
für jeden Punkt
\mathl{(x,y) \in \R^2}{} außerhalb der Einheitskreisscheibe zwei Urbildpunkte, auf dem Einheitskreis einen Urbildpunkt und innerhalb der offenen Einheitskreisscheibe keinen Urbildpunkt besitzt. Man interpretiere dies geometrisch.

}
{} {}

\inputaufgabe
{}
{

Es sei

\mavergleichskette
{\vergleichskette
{ M
}
{ \subseteq }{ N
}
{  }{
}
{    }{
}
{   }{
}
}
{}{}{}
eine
\definitionsverweis {abgeschlossene Untermannigfaltigkeit}{}{}
der
\definitionsverweis {differenzierbaren Mannigfaltigkeit}{}{}
$N$. Zeige, dass

\mavergleichskette
{\vergleichskette
{ TM
}
{ \subseteq }{ TN
}
{  }{
}
{    }{
}
{   }{
}
}
{}{}{}
eine abgeschlossene Teilmenge ist.

}
{} {}

\inputaufgabe
{}
{

Man gebe ein Beispiel einer
\definitionsverweis {injektiven}{}{}
\definitionsverweis {differenzierbaren Abbildung}{}{}
\maabbdisp {\varphi} { M } { N
} {}
zwischen zwei
\definitionsverweis {differenzierbaren Mannigfaltigkeiten}{}{}
\mathkor {} {M} {und} {N} {}
derart, dass die zugehörige
\definitionsverweis {Tangentialabbildung}{}{}
\maabbdisp {T(\varphi)} {TM} {TN
} {}
nicht injektiv ist.

}
{} {}

\inputaufgabe
{}
{

Man gebe ein Beispiel einer
\definitionsverweis {surjektiven}{}{}
\definitionsverweis {differenzierbaren Abbildung}{}{}
\maabbdisp {\varphi} { M } { N
} {}
zwischen zwei
\definitionsverweis {differenzierbaren Mannigfaltigkeiten}{}{}
\mathkor {} {M} {und} {N} {}
derart, dass die zugehörige
\definitionsverweis {Tangentialabbildung}{}{}
\maabbdisp {T(\varphi)} {TM} {TN
} {}
nicht surjektiv ist.

}
{} {}

\inputaufgabegibtloesung
{}
{

Man gebe ein
\definitionsverweis {stetiges}{}{} \definitionsverweis {Vektorfeld}{}{}
auf $S^2$ an, das nur eine Nullstelle besitzt.

}
{} {}

\inputaufgabe
{}
{

Die
\definitionsverweis {Einheitssphäre}{}{}

\mathl{S^2}{} bewege sich in einer Sekunde vollständig
\zusatzklammer {vom Nordpol aus gesehen gegen den Uhrzeigersinn} {} {}
um die Polachse. Welche differenzierbare Kurve und welcher Tangentialvektor an einen Punkt
\mathl{P \in S^2}{} gehört zu dieser Bewegung? Beschreibe das zugehörige
\definitionsverweis {Vektorfeld}{}{}
\maabbdisp {} {S^2} {TS^2 \subseteq T\R^3 = \R^6
} {.}

}
{} {}

Es sei $M$ eine
\definitionsverweis {differenzierbare Mannigfaltigkeit}{}{}
mit dem
\definitionsverweis {Tangentialbündel}{}{}
\maabb {p} {TM} {M
} {}
und sei
\maabbdisp {\psi} { Z } { M
} {}
eine
\definitionsverweis {Abbildung}{}{,}
wobei $Z$ eine Menge bezeichnet. Eine Abbildung
\maabbdisp {F} { Z } { TM
} {}
heißt
\definitionswort {Vektorfeld längs}{}
$\psi$, wenn

\mavergleichskette
{\vergleichskette
{ p \circ F
}
{ = }{ \psi
}
{  }{
}
{    }{
}
{   }{
}
}
{}{}{}
gilt.

  \zwischenueberschrift{Aufgaben zum Abgeben}

\inputaufgabe
{8 (3+3+2)}
{

Es seien

\mavergleichskette
{\vergleichskette
{ m,n
}
{ \in }{\N_+
}
{  }{
}
{    }{
}
{   }{
}
}
{}{}{.}

a) Zeige, dass die Menge

\mavergleichskettedisp
{\vergleichskette
{ M
}
{ =} { \left\{ (x,y,u,v) \in \R^4  \mid  ux^m+vy^n = 1         \right\}
}
{ } {
}
{ } {
}
{ } {
}
}
{}{}{}
eine
\definitionsverweis {abgeschlossene Untermannigfaltigkeit}{}{}
des $\R^4$ ist.

b) Zeige, dass die
\definitionsverweis {Abbildung}{}{}
\maabbeledisp {\varphi} { M } {\R^2
} {(x,y,u,v)} {(x,y)
} {,}
\definitionsverweis {differenzierbar}{}{}
und in jedem Punkt

\mavergleichskette
{\vergleichskette
{ P
}
{ \in }{ M
}
{  }{
}
{    }{
}
{   }{
}
}
{}{}{}
\definitionsverweis {regulär}{}{}
ist.

c) Beschreibe die
\definitionsverweis {Fasern}{}{}
von $\varphi$.

}
{} {}

\inputaufgabe
{10 (2+3+5)}
{

Wir betrachten die
\definitionsverweis {Abbildung}{}{}
\maabbeledisp {\varphi} {\R} {\R^2
} { x } {(x^2,x^3)
} {.}

a) Zeige, dass diese Abbildung
\definitionsverweis {differenzierbar}{}{}
und
\definitionsverweis {injektiv}{}{}
ist.

b) Zeige, dass $\varphi$ nicht in jedem Punkt
\definitionsverweis {regulär}{}{}
ist.

c) Zeige, dass das
\definitionsverweis {Bild}{}{}
von $\varphi$
\definitionsverweis {abgeschlossen}{}{}
in $\R^2$ ist, aber keine
\definitionsverweis {abgeschlossene Untermannigfaltigkeit}{}{}
des $\R^2$ ist.

}
{} {}

\inputaufgabe
{4}
{

Zeige, dass es eine
\definitionsverweis {Homöomorphie}{}{} des
\definitionsverweis {Tangentialbündels}{}{}
\mathl{T_{S_1 }}{} der
$1$-\definitionsverweis {Sphäre}{}{} $S^1$ mit dem
\definitionsverweis {Produkt}{}{} $S^1 \times \R$ gibt.

}
{} {}

\inputaufgabe
{5}
{

Es sei $M$ eine
\definitionsverweis {differenzierbare Mannigfaltigkeit}{}{}
und
\maabbdisp {\pi} {TM} {M
} {}
das
\definitionsverweis {Tangentialbündel}{}{.}
Zeige, dass
\mathl{TM}{} selbst in natürlicher Weise eine
\definitionsverweis {topologische Mannigfaltigkeit}{}{}
ist.

}
{} {}

In der folgenden Aufgabe wird der Begriff eines $R$-Moduls verwendet
\zusatzklammer {das ist eine direkte Verallgemeinerung des Vektorraumsbegriffes} {} {.}

Es sei $R$ ein
\definitionsverweis {kommutativer Ring}{}{}
und

\mavergleichskette
{\vergleichskette
{ M
}
{ = }{ (M,+,0)
}
{  }{
}
{    }{
}
{   }{
}
}
{}{}{}
eine \stichwort {additiv} {} geschriebene
\definitionsverweis {kommutative Gruppe}{}{.}
Man nennt $M$ einen
\definitionswortpraemath {R}{ Modul }{,}
wenn eine Operation
\maabbeledisp {} {R \times M } { M
} {(r,v)} { rv = r\cdot v
} {,}
\zusatzklammer {\stichwort {Skalarmultiplikation} {} genannt} {} {}
festgelegt ist, die folgende Axiome erfüllt
\zusatzklammer {dabei seien

\mavergleichskettek
{\vergleichskettek
{ r,s
}
{ \in }{ R
}
{  }{
}
{    }{
}
{   }{
}
}
{}{}{}
und

\mavergleichskettek
{\vergleichskettek
{ u,v
}
{ \in }{ M
}
{  }{
}
{    }{
}
{   }{
}
}
{}{}{}
beliebig} {} {:}
\aufzaehlungvier{
\mavergleichskette
{\vergleichskette
{ r(su)
}
{ = }{ (rs) u
}
{  }{
}
{    }{
}
{   }{
}
}
{}{}{,}
}{
\mavergleichskette
{\vergleichskette
{ r(u+v)
}
{ = }{ (ru) + (rv)
}
{  }{
}
{    }{
}
{   }{
}
}
{}{}{,}
}{
\mavergleichskette
{\vergleichskette
{ (r+s)u
}
{ = }{ (ru)+ (su)
}
{  }{
}
{    }{
}
{   }{
}
}
{}{}{,}
}{
\mavergleichskette
{\vergleichskette
{ 1 u
}
{ = }{  u
}
{  }{
}
{    }{
}
{   }{
}
}
{}{}{.}
}

\inputaufgabe
{4}
{

Es sei $M$ eine
\definitionsverweis {differenzierbare Mannigfaltigkeit}{}{,}
sei

\mavergleichskette
{\vergleichskette
{ R
}
{ = }{ C^1(M,\R)
}
{  }{
}
{    }{
}
{   }{
}
}
{}{}{}
der
\definitionsverweis {Ring}{}{}
der
\definitionsverweis {differenzierbaren Funktionen}{}{}
auf $M$ und sei $F$ die Menge aller
\definitionsverweis {Vektorfelder}{}{}
auf $M$.

a) Definiere eine Addition auf $F$ derart, dass $F$ zu einer
\definitionsverweis {kommutativen Gruppe}{}{}
wird.

b) Definiere eine Skalarmultiplikation
\maabbeledisp {} {R \times F } { F
} {(f,s)} {fs
} {,}
derart, dass $F$ zu einem
$R$-\definitionsverweis {Modul}{}{}
wird.

}
{} {}

 [[Kategorie:Latexseite]]
